% Source: https://artofproblemsolving.com/community/c3h2026917_moldy_cheese_mock_mathcounts_problems_110
% Stable identity: aops-topic-2026917-q10
Let $\pi(n)$ be the number of primes at most $n$; we need $\pi(n)\ge\dfrac {n}{2}$.

For $n\le 10$:
\[
\begin{array}{c|cccccccccc}
n & 1 & 2 & 3 & 4 & 5 & 6 & 7 & 8 & 9 & 10\\\hline
\pi(n) & 0 & 1 & 2 & 2 & 3 & 3 & 4 & 4 & 4 & 4\\
\dfrac {n}{2} & 0.5 & 1 & 1.5 & 2 & 2.5 & 3 & 3.5 & 4 & 4.5 & 5
\end{array}
\]
so $n=2,3,4,5,6,7,8$ work and $n=1,9,10$ do not.

For $n\ge9$, the non-primes up to $n$ include $1$, the $\left\lfloor\dfrac {n}{2}\right\rfloor-1$ even numbers from $4$ to $n$, and $9$. That is at least $\left\lfloor\dfrac {n}{2}\right\rfloor+1>\dfrac {n}{2}$ non-primes, so fewer than half of the numbers are prime. Hence exactly
\[\boxed{7}\]
positive integers work.
